\documentclass[../../main.tex]{subfiles}
\begin{document}

\paragraph{【新課程】}
{\bf [解]}
時刻$0$で$OX$を$x$軸正方向とし，$OX$が$x$軸から$y$軸正方向へと回転する座標系をとる．このとき，時刻$t$での点$P$の座標$(X, Y)$は
\begin{align}
  \begin{cases}
    X = e^{2t} \cos t \\
    Y = e^{2t} \sin t
  \end{cases}
\end{align}
と表される．

\begin{figure}[h]
  \centering
  \begin{tikzpicture}[scale=1.2, >=stealth]
    \draw[->] (-0.3,0) -- (2.5,0) node[right] {$x$};
    \draw[->] (0,-0.3) -- (0,2.0) node[above] {$y$};
    \node[below left] at (0,0) {$O$};
    \coordinate (P) at (1.8, 1.2);
    \draw[thick] (0,0) -- (P) node[above right] {$P$};
    \fill (P) circle (1.5pt);
    \draw[->] (0.6,0) arc (0:33.69:0.6);
    \node at (0.8,0.2) {$t$};
  \end{tikzpicture}
  \caption{点$P$の運動と座標系}
\end{figure}

両辺を$t$で微分すると，
\begin{align}
  \begin{cases}
    \dfrac{dX}{dt} = e^{2t} (2\cos t - \sin t) \\
    \dfrac{dY}{dt} = e^{2t} (2\sin t + \cos t)
  \end{cases}
\end{align}
となる．$0 \le t \le 2\pi$の区間で点$P$が同じ位置を通過することはないので，求める道のり$L$は
\begin{align}
  L &= \int_0^{2\pi} \sqrt{\left(\dfrac{dX}{dt}\right)^2 + \left(\dfrac{dY}{dt}\right)^2} \, dt \nonumber\\
  &= \int_0^{2\pi} e^{2t} \sqrt{(2\cos t - \sin t)^2 + (2\sin t + \cos t)^2} \, dt \nonumber\\
  &= \int_0^{2\pi} \sqrt{5} e^{2t} \, dt \nonumber\\
  &= \dfrac{\sqrt{5}}{2} \Bigl[ e^{2t} \Bigr]_0^{2\pi} \nonumber\\
  &= \dfrac{\sqrt{5}}{2} \left( e^{4\pi} - 1 \right) \tag{答}
\end{align}
である．

\paragraph{【旧課程】}
題意の定積分を$S$とおくと，
\begin{align}
  S &= \int_0^1 (3ax+2)(2x+b) \, dx \nonumber\\
  &= \int_0^1 \left\{ 6ax^2 + (3ab+4)x + 2b \right\} \, dx \nonumber\\
  &= \left[ 2ax^3 + \dfrac{3ab+4}{2}x^2 + 2bx \right]_0^1 \nonumber\\
  &= 2a + \dfrac{3ab+4}{2} + 2b \nonumber\\
  &= 2(a+b) + \dfrac{3}{2}ab + 2 = 0 \label{eq:1}
\end{align}
が成り立つ．
ここで$p = a+b$，$q = ab$とおいて，$a$，$b$が存在するための条件を調べる．
まず\eqref{eq:1}を$p,q$で表すと
\begin{align}
  2p + \dfrac{3}{2}q + 2 = 0 \Longleftrightarrow q = -\dfrac{4}{3}(p+1) \label{eq:2}
\end{align}
となる．
$a$，$b$は$t$についての$2$次方程式
\begin{align}
  t^2 - pt + q = 0
\end{align}
の$2$実数解であるから，判別式を$D$とすると$D \ge 0$が必要十分である．
\begin{align}
  D \ge 0 &\Longleftrightarrow p^2 - 4q \ge 0 \nonumber\\
  &\Longleftrightarrow p^2 + \dfrac{16}{3}(p+1) \ge 0 \nonumber\\
  &\Longleftrightarrow 3p^2 + 16p + 16 \ge 0 \nonumber\\
  &\Longleftrightarrow (3p+4)(p+4) \ge 0 \nonumber\\
  &\Longleftrightarrow p \le -4, \quad \dfrac{4}{3} \le p
\end{align}
となる．したがって，求める$a+b$の範囲は
\begin{align}
  a+b \le -4, \quad \dfrac{4}{3} \le a+b \tag{答}
\end{align}
である．

\end{document}
